\section{Model and exact sign program}

\begin{definition}[Admissible scores]\label{def:model}
Fix a finite discrete clock and distinct deterministic supports
$I_i=(a_i,b_i]$ with $a_i<b_i$. On a filtered probability space,
a score $X_i$ satisfies
\[
 |X_i|\le 1,\qquad X_i\text{ is }\mathcal F_{b_i}\text{-measurable},
 \qquad \E[X_i\mid\mathcal F_{a_i}]=0.
\]
All scores maturing at a given time are available at issues at that time.
The weight vector $w$ is real and deterministic. A family is laminar
when any two supports are disjoint or one contains the other. Its depth
is the number of supports in a longest strict inclusion chain.
\end{definition}

The admissible class varies over filtered spaces and laws. All scores
have mean zero, so define
\begin{align}
 \Vstar &= \sup_{\rm admissible}\E\Big(\sum_i w_i X_i\Big)^2,\\
 U &= \max\{w^\top Cw : C=C^\top\succeq0,\ C_{ii}=1,\
           C_{ij}=0\text{ if } I_i\cap I_j=\varnothing\}.\label{eq:sdp}
\end{align}
Here $\Vstar$ is the worst variance for this law class. Replace $X_i$ by
$\operatorname{sgn}(w_i)X_i$ with either sign at a zero weight, and
conjugate $C$ by the same diagonal sign matrix. This preserves both
problems and permits nonnegative weights below. Other deterministic
score bounds can be absorbed into the weights. The identity is feasible
in~\eqref{eq:sdp}, so $U>0$ when $w\ne0$. All zero-weight instances
have both optima zero and require no ratio.

\begin{theorem}[Full matured-history sign LP]\label{thm:lp}
Let $M_i=\{j:b_j\le a_i\}$. The value $\Vstar$ equals the maximum of
$\sum_z p_z(w^\top z)^2$ over the following finite polytope:
\[
 p_z\ge0,\qquad \sum_{z\in\{-1,1\}^n}p_z=1,\qquad
 \sum_{z:z_{M_i}=h}p_z z_i=0
 \quad\bigl(i,\ h\in\{-1,1\}^{M_i}\bigr).
\]
For each convex function of the weighted total, maximizing its
expectation over this polytope gives the corresponding supremum.
\end{theorem}
\begin{proof}
On an independent product extension sample independent uniform
random variables $L_i$ on $[0,1]$. Release $L_i$ at $b_i$ and set
$Z_i=2\mathbf1_{\{L_i\le(1+X_i)/2\}}-1$.
Conditional on the original terminal field the signs are independent
with means $X_i$. Conditional Jensen therefore increases every convex
payoff, and the square increases in expectation by
$\sum_i w_i^2(1-\E X_i^2)$. At $a_i$ the variable $L_i$ is fresh.
Condition first on the original terminal field and the released uniforms
to integrate $L_i$, then project down to the original issue field
enlarged by the released uniforms. The conditional expectation of $X_i$
there equals $\E[X_i\mid\mathcal F_{a_i}]=0$ by independence. Projecting
down to the matured signs gives all the displayed equations.

Conversely, sample a sign vector with masses $p_z$ and at time $t$ take
$\sigma(Z_j:b_j\le t)$ as the filtration. Each positive-mass history
has conditional mean zero by its equation, while a zero-mass history
imposes no further condition. This realizes every feasible vector of
masses. Compactness of the nonempty polytope gives attainment.
\end{proof}

If $b_i\le a_j$, the tower property gives $\E X_iX_j=0$. Thus the sign
problem is bounded above by~\eqref{eq:sdp}. Pairwise zero covariances alone
do not enforce conditional centering. For three disjoint consecutive
intervals, let the signs be $R,S,RS$ for independent fair signs $R,S$.
The covariance matrix is the identity, but $RS$ is already determined
at its issue. In the opposite direction, admissible disjoint subtree
vectors may be dependent while each cross covariance vanishes.

\section{Laminar convex-order and Gram recursions}

For bounded random variables $Y,Y'$ write $Y\cx Y'$ when
$\E f(Y)\le \E f(Y')$ for every real convex $f$. Write $\sgn_0$ for
the sign with an independent fair tie coin at zero.

\begin{lemma}[Fair root extension]\label{lem:root}
If $S\cx S^\circ$, $S^\circ$ is symmetric, $Y$ is a fair sign with
arbitrary coupling to $S$, and $m\ge0$, then
\[
 S+mY\cx S^\circ+m\sgn_0(S^\circ).
\]
\end{lemma}
\begin{proof}
Fix convex $f$ and put $\lambda=(f(m)-f(-m))/2$ and
$g(s)=\max\{f(s+m)-\lambda,f(s-m)+\lambda\}$.
The convex function $g$ bounds $f(S+mY)-\lambda Y$ pointwise.
The expectations are bounded by $\E g(S^\circ)$ because $Y$ is fair
and $g$ is convex. The increment $f(s+m)-f(s-m)$ is nondecreasing
in $s$, and the two branches meet at zero. Hence $\sgn_0(S^\circ)$
selects a maximizing branch and is fair by symmetry. Substituting
this choice shows that $\E g(S^\circ)$ equals the stated upper payoff.
The argument also covers a zero weight and atoms at zero.
\end{proof}

The selected coupling is comonotone. More generally, with a fixed
bounded inner law $S$, the largest $\E[SY]$ over fair signs $Y$ is
$\min_t\E|S-t|$. For any $t$, fairness gives
$\E[SY]=\E[(S-t)Y]\le \E|S-t|$. At a median use the sign of $S-t$
off the median and choose the tie probabilities to make the sign fair.
If a root has a trivial incoming field and can observe the inner vector
at maturity, this coupling is admissible. The threshold zero in
Lemma~\ref{lem:root} is justified by symmetry of the reference law.

\begin{theorem}[Conditional laminar extremum]\label{thm:tree}
Order the children of every node $v$ of a finite laminar forest by time.
Independently sample their canonical vectors and a fresh fair tie coin
at $v$. Recursively define the canonical total
\[
 S_v=\sum_{c:\,c\text{ child of }v}T_c,\qquad
 T_v=S_v+w_v\sgn_0(S_v).
 \tag{*}\label{eq:tree}
\]
The empty sum is zero. For every admissible law and event
$E\in\mathcal F_{a_v}$ with $\Pr(E)>0$, the conditional distribution
of the weighted subtree total given $E$ is dominated by $T_v$ in
convex order. The canonical vector is admissible in its matured-sign
filtration. Independent maximal components give a convex-order
maximum for the forest total and
\[
 \Vstar=\sum_{v\text{ maximal}}\E T_v^2.
\]
\end{theorem}
\begin{proof}
Prove the eventwise statement by induction on the subtree. For bounded
scores, endpoint rounding as in Theorem~\ref{thm:lp} also applies
under the law conditioned on $E$. In particular, $E$ belongs to the
original terminal field, so conditional Jensen given that field and
averaging on $E$ yield the convex comparison with the rounded signs.
It suffices to prove the induction for signs. Conditioning on $E$
preserves every descendant's issue constraint because $E$ belongs to
each descendant's issue field. For any event $F$ in that issue field,
the original centering equation integrated on $E\cap F$ gives the
conditional-law equation on $F$.
A leaf is fair under this conditioning and has the canonical law.

Under the law conditioned on $E$, fix the complete sign vectors of
the children preceding the last child. Their maturities are no later
than the last child's issue. Each positive-mass joint history, together
with $E$, is an incoming event for that subtree. On this history the
earlier total is a fixed number. The induction hypothesis applied to
the convex payoff translated by that total bounds the expected payoff
by averaging against the same canonical last-child law for every
history. Such averaging leaves a convex function of the earlier sum.
Iterating backwards over the children replaces the adaptive joint sum
by the sum of independent canonical child totals in convex order.
This reference sum is symmetric. The actual root sign is fair on $E$.
Lemma~\ref{lem:root} proves the eventwise comparison with~\eqref{eq:tree}.

For attainment generate the canonical signs from independent primitive
tie coins at the nodes. A node's sign is a function of its own subtree's
coins and is fair by symmetry, including at zero. At its issue every
descendant is still immature because each has positive length and starts
at or after that issue. Every ancestor is also immature. The matured
scores are in separated branches and are functions of coins
outside the node's subtree. Thus the sign is conditionally fair at issue. The
matured-sign filtration reveals all coincident maturities together and
makes the assigned signs measurable at maturity. This establishes
simultaneous feasibility. Apply the same backwards argument to the
ordered maximal components, without an additional root, to obtain the
forest convex-order maximum. Independent component totals are centered,
so the square yields the formula for $\Vstar$.
\end{proof}

Independent child vectors are one attaining choice. For example, with
two disjoint two-leaf stars of unit weights inside one root, the later
root may copy its first leaf if the earlier first leaf is positive and
its second leaf otherwise. Its conditional total has the same law on
both histories, while its joint vector depends on the first star.
The outer sign can still attain the optimum.

\begin{proposition}[Gram recursion]\label{prop:gram}
Define $u_v=w_v+\sqrt{\sum_{c:\,c\text{ child of }v}u_c^2}$. Then
$U=\sum_{v\text{ maximal}}u_v^2$.
\end{proposition}
\begin{proof}
Represent $C$ by unit Gram vectors. The spans of distinct child
subtrees are orthogonal because their supports are disjoint. The
largest possible child-sum norm is bounded by the square root of the
sum of their squared optima. Adding the parent increases it by at most
$w_v$. To attain both bounds, embed optimal child systems in orthogonal
spaces and align the parent with their weighted sum. If it vanishes,
use a fresh orthogonal unit vector. No disjointness zero inside the
subtree involves its parent. At a maximal component the same construction
leaves its span orthogonal to every other maximal component.
\end{proof}
